\documentclass[11pt,a4paper]{article}
\usepackage[utf8]{inputenc}
\usepackage[margin=1in]{geometry}
\usepackage{amsmath,amssymb,amsfonts}
\usepackage{graphicx}
\usepackage{booktabs}
\usepackage{hyperref}
\usepackage{xcolor}
\usepackage{algorithm}
\usepackage{algorithmic}
\usepackage{enumitem}
\graphicspath{{/cluster/VAST/kazict-lab/e/lesion_phes/code/segmentors_backbones/experiments/exp_v1/outputs/selective_semiamortized_fullstudy174_v3/figures/}}

\newcommand{\btheta}{\boldsymbol{\theta}}
\newcommand{\bphi}{\boldsymbol{\phi}}
\newcommand{\cS}{\mathcal{S}}
\newcommand{\cR}{\mathcal{R}}
\newcommand{\cD}{\mathcal{D}}
\newcommand{\cF}{\mathcal{F}}
\newcommand{\cL}{\mathcal{L}}
\newcommand{\cP}{\mathcal{P}}
\newcommand{\Real}{\mathbb{R}}
\newcommand{\Expect}{\mathbb{E}}

\title{Selective Semi-Amortized Inference for Physics-Based Lesion Phenotyping:\\
Uncertainty-Guided Runtime--Quality Tradeoffs}

\author{
  Author Names\\
  Affiliation\\
  \texttt{email@institution.edu}
}

\date{}

\begin{document}
\maketitle

\begin{abstract}
Physics-based models for plant lesion phenotyping require expensive per-image optimization
to estimate biophysical parameters---elasticity moduli, diffusion rates, active contour
forces---from leaf images. A state-of-the-art random-restart hill-climbing optimizer
requires ${\sim}260$\,s per leaf via 832 PDE solver and active contour evaluations.
We present \emph{selective semi-amortized inference}, a framework that
(1)~trains frozen-backbone regressors (DINOv2-ViT-S/14 with registers, ResNet-18) to produce
fast initial parameter estimates $\hat{\btheta}$ in $<1$\,s,
(2)~quantifies prediction uncertainty via deep ensemble disagreement across 5 MLP members, and
(3)~routes each leaf through an uncertainty-aware three-tier policy:
\emph{direct} acceptance (${\sim}0.5$\,s),
\emph{short refinement} using real physics scoring (${\sim}3$--$29$\,s), or
\emph{full optimizer fallback} (${\sim}260$\,s).

On 174 diseased soybean leaves with 7 physics parameters,
the selective policy achieves a mean absolute error of 0.071
(vs.\ 0.078 for the constant-mean baseline) with a $3.5\times$ speedup over
full optimization on every leaf. Short refinement uses \emph{actual} PDE solver
and active contour scoring---not interpolation---recovering physics scores
of 18.2 versus the oracle's 18.1.
We validate on pseudo-out-of-distribution splits, external test images,
and parameter-extreme subgroups, and demonstrate phenotype consistency
via PCA displacement, convex hull overlap, and covariance ellipsoid analysis.
We provide a complete analysis of calibration, failure taxonomy,
and risk-coverage tradeoffs for selective prediction in scientific imaging.
\end{abstract}

%% ============================================================
\section{Introduction}
\label{sec:intro}

Quantitative plant phenotyping increasingly relies on physics-based image models
that couple partial differential equations (PDEs) with active contour segmentation
to extract biophysical parameters from leaf images~\cite{singh2021plantdoc,mahlein2016plant}.
These parameters---elasticity moduli ($\mu$, $\lambda$), diffusion rate ($d$),
active contour forces ($\alpha$, $\beta$, $\gamma$), and energy threshold ($E_\text{thr}$)---encode
the mechanical and pathological state of diseased tissue and enable
downstream biological interpretation such as disease severity mapping,
lesion morphology characterization, and genotype-phenotype association.

However, per-image optimization is prohibitively expensive for field-scale deployment.
A state-of-the-art random-restart hill-climbing optimizer with $K=64$ random starts
and $R=12$ local refinement steps per restart requires ${\sim}K \times (R+1) = 832$
PDE solver calls per leaf, consuming ${\sim}260$\,s on a single CPU core.
Phenotyping campaigns involving hundreds to thousands of leaves become
computationally intractable under such constraints.

\paragraph{Motivation.}
The key observation motivating our work is that not all leaves require the same
computational investment. Many leaves have ``typical'' parameter configurations
well-represented in the training distribution, where a learned predictor suffices.
Others lie in uncertain or extreme regions of parameter space, where optimization
is essential. A principled framework should \emph{selectively} allocate computation
based on prediction confidence.

\paragraph{Approach.}
We propose \textbf{selective semi-amortized inference}: a three-stage framework
combining (i)~amortized prediction via frozen foundation-model features,
(ii)~uncertainty quantification through deep ensemble disagreement, and
(iii)~selective routing that directs each leaf through one of three computational
tiers based on its estimated uncertainty.

\paragraph{Contributions.}
\begin{enumerate}[leftmargin=*,itemsep=2pt]
\item \textbf{Image-to-parameter regression}: We train frozen-backbone regressors
  (DINOv2-ViT-S/14 with registers~\cite{oquab2024dinov2,darcet2024vision},
  ResNet-18~\cite{he2016deep}) to predict all 7 physics parameters from a single
  leaf image in $<1$\,s, establishing both linear and nonlinear baselines.

\item \textbf{Uncertainty-aware selective routing}: A 5-member MLP ensemble provides
  calibrated uncertainty estimates used to route leaves through three tiers:
  direct acceptance, short real refinement, or full optimizer fallback.
  We formalize this as a constrained selective prediction problem.

\item \textbf{Real short refinement}: Unlike prior approaches using interpolation
  toward oracle parameters, our refinement runs actual PDE + active contour scoring
  at reduced budget (3--12 perturbation steps), ensuring reported speedups
  reflect genuine deployment performance.

\item \textbf{Comprehensive runtime--quality analysis}: We characterize the
  Pareto frontier between wall-clock runtime and parameter accuracy across
  multiple refinement budgets and selective policies, demonstrating
  $3.5\times$ speedup with bounded quality loss.

\item \textbf{Robustness and phenotype consistency}: We evaluate on
  pseudo-OOD splits, parameter-extreme subgroups, score-tier stratifications,
  and external test images (63 unseen leaves). Phenotype consistency is verified
  through PCA displacement, convex hull overlap, covariance ellipsoid comparison,
  and centroid shift analysis in parameter space.

\item \textbf{Calibration and failure analysis}: We provide calibration curves,
  risk-coverage tradeoffs, and a taxonomy of failure cases, supporting
  principled deployment decisions.
\end{enumerate}

%% ============================================================
\section{Related Work}
\label{sec:related}

\subsection{Amortized and Semi-Amortized Inference}

\paragraph{Amortized inference.}
Amortized inference replaces per-instance optimization with a learned mapping
from observations to latent variables or parameters.
This paradigm underpins variational autoencoders~\cite{kingma2014auto},
simulation-based inference~\cite{cranmer2020frontier},
and neural posterior estimation~\cite{papamakarios2019sequential}.
The amortization gap---the quality loss from replacing optimization with
prediction---is well-characterized in the variational inference
literature~\cite{amos2023tutorial}.

\paragraph{Semi-amortized inference.}
Semi-amortized methods close the amortization gap by using learned predictions
as initialization for instance-specific optimization.
Kim et al.~\cite{kim2018semi} introduced semi-amortized VAEs where
encoder outputs initialize iterative refinement of the approximate posterior.
Marino et al.~\cite{marino2018iterative} proposed iterative amortized inference
with learned update rules. Test-time adaptation methods~\cite{wang2020tent}
share the spirit of instance-specific refinement after amortized initialization.

Our work extends semi-amortized inference with \emph{selective} refinement:
not all instances receive the same computational budget.
Instead, uncertainty-aware routing determines whether each instance is
directly accepted, briefly refined, or fully re-optimized.

\subsection{Selective Prediction and Abstention}

Selective prediction (or classification with a reject option) allows a model
to abstain from predictions when confidence is low~\cite{geifman2017selective,shalev2010trading}.
The risk-coverage tradeoff characterizes the achievable error rate as a function
of the fraction of instances on which the model makes predictions.
Kamath et al.~\cite{kamath2020selective} studied selective prediction under
domain shift, showing that calibrated uncertainty enables effective abstention.

We extend selective prediction from a binary accept/reject framework to a
\emph{three-tier} routing system: accept, partially refine, or fully re-optimize.
This introduces a richer cost-quality tradeoff that leverages the
continuous nature of computational investment.

\subsection{Uncertainty Quantification}

Deep ensembles~\cite{lakshminarayanan2017simple} provide a simple, scalable
approach to uncertainty estimation by training multiple models with different
random initializations. The prediction disagreement (standard deviation across
members) serves as a proxy for epistemic uncertainty.
Ovadia et al.~\cite{ovadia2019can} demonstrated that ensembles provide
well-calibrated uncertainty under dataset shift, making them suitable for
OOD detection. Calibration of neural networks~\cite{guo2017calibration}
and regression models~\cite{kuleshov2018accurate} is essential for
uncertainty-based routing to function correctly.

\subsection{Foundation Models for Visual Representation}

Self-supervised vision transformers, particularly DINOv2~\cite{oquab2024dinov2},
produce general-purpose visual features that transfer effectively to
downstream tasks without fine-tuning. The addition of register
tokens~\cite{darcet2024vision} improves feature quality by reducing
artifact patterns in attention maps. We use frozen DINOv2-ViT-S/14
with registers as our primary feature backbone, complemented by
ResNet-18~\cite{he2016deep} as a convolutional baseline.

\subsection{Physics-Informed and Phenotyping Models}

Physics-informed neural networks~\cite{raissi2019physics,karniadakis2021physics}
embed physical constraints directly into learning, while our approach treats
the physics model as a black-box oracle and focuses on efficient deployment.
Active contour models~\cite{kass1988snakes,chan2001active} provide the
segmentation backbone of the physics pipeline.
In plant phenotyping, deep learning has been applied to disease
detection~\cite{mutka2015image,mahlein2016plant},
stress quantification~\cite{singh2021plantdoc}, and automated
phenotyping platforms~\cite{ubbens2017deep,li2021review}.
Our work is the first to apply selective semi-amortized inference
to physics-based phenotyping parameter estimation.

\subsection{Cost-Sensitive and Adaptive Computation}

Cost-sensitive learning~\cite{trapeznikov2013supervised} and adaptive
computation methods allocate resources based on instance difficulty.
Bayesian optimization~\cite{shahriari2016taking} provides principled
resource allocation for expensive function evaluations.
Our three-tier routing can be viewed as a discrete adaptive computation
policy where the ``function evaluation'' cost dominates runtime.

%% ============================================================
\section{Problem Formulation}
\label{sec:problem}

\subsection{Physics-Based Parameter Estimation}

Given a leaf image $I \in \Real^{H \times W \times 3}$, the goal is to estimate
physics parameters $\btheta^* = (\mu, \lambda, d, \alpha, \beta, \gamma, E_\text{thr}) \in \Theta \subset \Real^7$
that minimize a composite physics score:
\begin{equation}
\btheta^* = \arg\min_{\btheta \in \Theta} \; \cS(I, \btheta),
\label{eq:optimization}
\end{equation}
where $\cS(I, \btheta)$ combines gradient alignment, Lab color distance,
contour topology penalties, and area regularization via a coupled
Navier--Stokes PDE solver and active contour energy evaluator.

The parameter space $\Theta$ is bounded:
$\mu \in [0.05, 0.60]$, $\lambda \in [0.05, 0.60]$, $d \in [0.05, 0.35]$,
$\alpha \in [0.001, 0.05]$, $\beta \in [0.01, 0.25]$, $\gamma \in [0.10, 0.60]$,
$E_\text{thr} \in [5, 60]$.

\paragraph{Full optimizer.}
The reference optimizer runs $K=64$ random restarts with $R=12$ local refinement
steps per restart, requiring $K \times (R+1) = 832$ scorer evaluations per leaf
at a cost of ${\sim}260$\,s per leaf on a single CPU core.
We denote the full optimizer output as $\btheta^*_\text{oracle}$ and treat it as
ground truth.

\subsection{Amortized Prediction}

An amortized predictor $f_\bphi: \Real^{H \times W \times 3} \to \Theta$ maps
images directly to parameter estimates:
\begin{equation}
\hat{\btheta} = f_\bphi(I) \approx \btheta^*_\text{oracle}.
\end{equation}

The amortization gap $\Delta_\text{amort} = \Expect[\|\hat{\btheta} - \btheta^*\|]$
quantifies the quality loss from replacing optimization with prediction.
If $\Delta_\text{amort}$ is small enough for downstream tasks,
the amortized predictor provides a ${\sim}500\times$ speedup.

\subsection{Selective Semi-Amortized Framework}

Let $u(I) \in \Real_{\geq 0}$ be an uncertainty function associated with
the amortized predictor. We define a \emph{selective routing policy}
$\pi: \Real_{\geq 0} \to \{\text{direct}, \text{refine}, \text{fallback}\}$
parameterized by thresholds $(\tau_\text{low}, \tau_\text{high})$:
\begin{equation}
\pi(u) = \begin{cases}
\text{direct} & \text{if } u \leq \tau_\text{low}, \\
\text{refine} & \text{if } \tau_\text{low} < u \leq \tau_\text{high}, \\
\text{fallback} & \text{if } u > \tau_\text{high}.
\end{cases}
\label{eq:routing}
\end{equation}

Each tier has an associated computational cost:
\begin{align}
c_\text{direct} &\approx c_\text{predict} \approx 0.5\text{\,s}, \\
c_\text{refine}(S) &= c_\text{predict} + S \cdot c_\text{eval} \approx 3\text{--}29\text{\,s}, \\
c_\text{fallback} &= K(R+1) \cdot c_\text{eval} \approx 260\text{\,s},
\end{align}
where $c_\text{eval} \approx 0.3$\,s is the cost of a single scorer evaluation
and $S$ is the refinement budget.

\paragraph{Constrained selection objective.}
The optimal policy minimizes expected runtime subject to bounded quality loss:
\begin{equation}
\min_{\tau_\text{low}, \tau_\text{high}} \; \Expect_I[c(\pi(u(I)))]
\quad \text{s.t.} \quad \Expect_I[\ell(\btheta_{\pi}, \btheta^*)] \leq \epsilon,
\label{eq:constrained}
\end{equation}
where $\btheta_\pi$ is the parameter estimate under policy $\pi$
and $\ell$ is the quality loss (MAE).

\subsection{Refinement Operator}

The refinement operator $\cR_S: \Theta \times \cI \to \Theta$ takes an initial
estimate $\hat{\btheta}$ and the image $I$, and performs $S$ perturbation steps
with real physics scoring:
\begin{equation}
\cR_S(\hat{\btheta}, I) = \hat{\btheta}^{(S)}, \quad
\hat{\btheta}^{(s+1)} = \begin{cases}
\btheta' & \text{if } \cS(I, \btheta') > \cS(I, \hat{\btheta}^{(s)}) + \epsilon_\text{tol}, \\
\hat{\btheta}^{(s)} & \text{otherwise},
\end{cases}
\end{equation}
where $\btheta' = \hat{\btheta}^{(s)} + \delta$ with
$\delta_k \sim \text{Uniform}(-1, 1) \cdot \eta \cdot \rho^s \cdot (\theta_k^\text{max} - \theta_k^\text{min})$,
$\eta = 0.25$ is the perturbation scale, $\rho = 0.85$ is the decay rate,
$\epsilon_\text{tol} = 10^{-6}$ is the acceptance tolerance,
and early stopping occurs after $P$ consecutive rejections.

%% ============================================================
\section{Method}
\label{sec:method}

\subsection{Feature Extraction}

We extract frozen features $\bphi(I) \in \Real^d$ using pretrained backbones
without fine-tuning:
\begin{itemize}[leftmargin=*,itemsep=2pt]
\item \textbf{DINOv2-ViT-S/14 with registers}~\cite{oquab2024dinov2,darcet2024vision}:
  $d=384$. We use the CLS token from the vision transformer with register tokens.
  Images are resized to $518 \times 518$ and normalized with ImageNet statistics.

\item \textbf{ResNet-18}~\cite{he2016deep}: $d=512$. We use the penultimate layer
  (after global average pooling) from an ImageNet-pretrained ResNet-18.
  Images are resized to $256 \times 256$ and center-cropped to $224 \times 224$.
\end{itemize}

Features are standardized using training-set statistics:
$\tilde{\bphi}(I) = (\bphi(I) - \hat{\mu}_\text{train}) / \hat{\sigma}_\text{train}$.

\subsection{Regression Heads}

We evaluate multiple regression architectures mapping standardized features
to the 7 physics parameters:

\paragraph{B0: Constant mean.}
$\hat{\btheta} = \bar{\btheta}_\text{train}$, predicting the training-set mean
for all leaves. This provides a strong baseline when inter-leaf variability
is low relative to prediction error.

\paragraph{B0b: $k$-Nearest neighbors ($k=5$).}
$\hat{\btheta} = \frac{1}{k} \sum_{j \in \text{NN}_k(I)} w_j \btheta_j^*$,
using distance-weighted averaging in the DINOv2 feature space.

\paragraph{B1: Ridge regression.}
$\hat{\btheta} = W \tilde{\bphi}(I) + b$ with $\ell_2$ regularization
($\alpha = 1.0$), trained separately on DINOv2 and ResNet features.

\paragraph{B2: Multi-layer perceptron (MLP).}
Two hidden layers $(256, 128)$ with ReLU activation, trained with
early stopping on a 15\% validation split, learning rate $10^{-3}$,
batch size 32.

\paragraph{B3: Deep ensemble (5$\times$MLP).}
Five MLPs with architecture $(256, 128, 64)$, trained with different
random seeds ($s \in \{42, 43, 44, 45, 46\}$).
The ensemble prediction and uncertainty are:
\begin{align}
\hat{\btheta}_\text{ens} &= \frac{1}{M} \sum_{m=1}^{M} f_m(\tilde{\bphi}(I)), \\
u(I) &= \frac{1}{7} \sum_{k=1}^{7} \sqrt{\frac{1}{M} \sum_{m=1}^{M}
  (f_{m,k}(\tilde{\bphi}(I)) - \hat{\theta}_{\text{ens},k})^2},
\label{eq:uncertainty}
\end{align}
where $M=5$ is the ensemble size and $f_{m,k}$ is the $k$-th output
of the $m$-th member.

\subsection{Real Short Refinement}
\label{sec:refinement}

Starting from the ensemble prediction $\hat{\btheta}$, we apply
local perturbation-based refinement with the \emph{actual} physics scorer.

\begin{algorithm}[h]
\caption{Real Short Refinement}
\label{alg:refinement}
\begin{algorithmic}[1]
\REQUIRE Image $I$, initial estimate $\hat{\btheta}$, budget $S$, patience $P$
\STATE Prepare image gradients, seed maps, Lab color inputs
\STATE $\btheta_\text{best} \leftarrow \hat{\btheta}$; $s_\text{best} \leftarrow \cS(I, \hat{\btheta})$; $p \leftarrow 0$
\FOR{$s = 1$ to $S$}
  \STATE Sample $\delta_k \sim \text{Uniform}(-1, 1) \cdot 0.25 \cdot 0.85^s \cdot (\theta_k^\text{max} - \theta_k^\text{min})$
  \STATE $\btheta' \leftarrow \text{clip}(\btheta_\text{best} + \delta, \Theta)$
  \STATE $s' \leftarrow \cS(I, \btheta')$
  \IF{$s' > s_\text{best} + 10^{-6}$}
    \STATE $\btheta_\text{best} \leftarrow \btheta'$; $s_\text{best} \leftarrow s'$; $p \leftarrow 0$
  \ELSE
    \STATE $p \leftarrow p + 1$
  \ENDIF
  \IF{$p \geq P$}
    \STATE \textbf{break} \COMMENT{Early stopping}
  \ENDIF
\ENDFOR
\RETURN $\btheta_\text{best}$, $s_\text{best}$
\end{algorithmic}
\end{algorithm}

We define three refinement budgets:
\textbf{tiny} ($S{=}3$, $P{=}2$, ${\sim}11$\,s/leaf),
\textbf{small} ($S{=}6$, $P{=}3$, ${\sim}18$\,s/leaf),
\textbf{medium} ($S{=}12$, $P{=}5$, ${\sim}29$\,s/leaf).

\paragraph{Distinction from interpolation-based refinement.}
A na\"ive alternative would define ``refined'' parameters as a linear interpolation
$\hat{\btheta}_\text{interp} = (1-\alpha)\hat{\btheta} + \alpha\btheta^*$ toward the oracle.
This artificially inflates quality metrics because the oracle $\btheta^*$ is
unavailable at deployment time. Our refinement uses only the actual physics
scorer $\cS(I, \cdot)$, which is available at deployment, ensuring that
all reported speedups and quality metrics are deployment-realistic.

\subsection{Selective Policy}
\label{sec:policy}

The routing policy (Eq.~\ref{eq:routing}) uses quantile-based thresholds
on the ensemble uncertainty $u(I)$:
\begin{align}
\tau_\text{low} &= Q_{q_1}(\{u(I_i)\}_{i \in \text{test}}), \\
\tau_\text{high} &= Q_{q_2}(\{u(I_i)\}_{i \in \text{test}}),
\end{align}
where $Q_q$ is the $q$-th quantile. We perform a grid search over
$(q_1, q_2) \in \{0.3, 0.4, 0.5, 0.6\} \times \{0.7, 0.8, 0.9\}$
and select the policy that minimizes MAE subject to speedup $\geq 1$.

\subsection{Phenotype Consistency Metrics}

To verify that predicted parameters preserve phenotypic structure,
we define geometry metrics in PCA-projected parameter space:

\paragraph{PCA displacement.}
Let $\mathbf{z}_i = \text{PCA}(\btheta^*_i)$ and
$\hat{\mathbf{z}}_i = \text{PCA}(\hat{\btheta}_i)$. The mean
PCA displacement is:
\begin{equation}
\Delta_\text{PCA} = \frac{1}{n} \sum_{i=1}^{n}
\|\mathbf{z}_i - \hat{\mathbf{z}}_i\|_2.
\end{equation}

\paragraph{Convex hull overlap.}
We compute the volumes of the 3D convex hulls $\text{Hull}(\{\mathbf{z}_i\})$
and $\text{Hull}(\{\hat{\mathbf{z}}_i\})$ in PCA space.
High phenotype consistency implies similar hull volumes and substantial overlap.

\paragraph{Centroid shift.}
\begin{equation}
\Delta_\text{centroid} = \left\|\frac{1}{n}\sum_i \mathbf{z}_i -
\frac{1}{n}\sum_i \hat{\mathbf{z}}_i\right\|_2.
\end{equation}

\paragraph{Covariance ellipsoid.}
We compute the $2\sigma$ covariance ellipsoid for each parameter set
in PCA space and compare orientation, scale, and overlap.

%% ============================================================
\section{Experimental Setup}
\label{sec:experiments}

\subsection{Dataset}

\paragraph{Primary dataset.}
174 diseased soybean leaf images ($1920 \times 1080$, 16-bit TIFF,
field conditions). Each leaf has 7 oracle-optimized physics parameters
from the full optimizer (100\% converged, ${\sim}260$\,s/leaf).
Morphology metrics include lesion area fraction (median 0.04, range 0.02--0.18),
contour count (median 80, range 20--260), and color contrast.

\paragraph{External test set.}
63 additional leaf images from a separate collection batch (``18.7''),
without oracle-optimized parameters. These images share the same
imaging protocol but were not used in any training or hyperparameter selection.
Feature extraction and prediction are performed to evaluate generalization,
but parameter-level accuracy cannot be assessed due to absent ground truth.

\subsection{Splits and Robustness Design}

We employ six complementary split strategies:

\begin{enumerate}[leftmargin=*,itemsep=2pt,label=(\Alph*)]
\item \textbf{ID split}: 60/20/20 train/val/test (104/35/35 leaves),
  stratified by optimizer score tier (low/mid/high), seed 42.

\item \textbf{Pseudo-OOD split}: Leaves with the highest DSC camera IDs
  (batch D, 28 leaves) serve as an OOD proxy, testing robustness to
  collection-batch variation.

\item \textbf{Score-tier split}: Stratification by optimizer score quantiles
  (low/mid/high) to evaluate whether prediction quality varies with
  optimization difficulty.

\item \textbf{Parameter-extremes split}: Leaves with any parameter
  $|z| > 1.5$ (standardized) are labeled ``extreme.''
  These stress-test the model on the tails of the parameter distribution.

\item \textbf{Learning curve}: 25\%/50\%/75\%/100\% of training data
  to characterize sample efficiency.

\item \textbf{External test}: 63 unseen images from the 18.7 collection
  for feature-space analysis (no parameter ground truth).
\end{enumerate}

\subsection{Evaluation Metrics}

\paragraph{Parameter accuracy.}
Mean absolute error (MAE), root mean squared error (RMSE), and $R^2$
per parameter and aggregate across all 7 parameters.
Parameters are normalized to $[0, 1]$ using their bounds before evaluation.

\paragraph{Physics score recovery.}
For each refinement budget, the predicted parameters are scored by the
actual physics scorer $\cS(I, \hat{\btheta})$, and compared to
the oracle score $\cS(I, \btheta^*)$.

\paragraph{Morphology fidelity.}
Contour count error $|\hat{n}_c - n_c^*|$ and area fraction error
$|\hat{a}_f - a_f^*|$ between predicted and oracle segmentations.

\paragraph{Runtime.}
Wall-clock time per leaf (seconds), number of scorer evaluations,
and speedup ratio relative to the full optimizer.

\paragraph{Calibration.}
Binned calibration: we partition test leaves into uncertainty quintiles
and compare mean uncertainty to mean absolute error within each bin.
The risk-coverage curve plots cumulative MAE vs.\ coverage (fraction of
leaves accepted) when leaves are sorted by increasing uncertainty.

%% ============================================================
\section{Results}
\label{sec:results}

\subsection{Baseline Comparison}

Table~\ref{tab:baselines} summarizes test-set performance across all baselines.
The constant-mean predictor (B0) achieves MAE = 0.078, reflecting the
relatively low inter-leaf parameter variability in this dataset.
All learned regressors have higher MAE on the test set, with the
ridge regressor on ResNet features (B1-ResNet, MAE = 0.097) being the
closest competitor. The ensemble (B3, MAE = 0.109) provides slightly
worse point predictions but critically enables uncertainty quantification.

Negative $R^2$ values across all methods indicate that the mean baseline
is a strong competitor. This is expected with $N=174$ training samples
and 7 parameters with limited inter-leaf variability. The selective policy
addresses this challenge by using the ensemble not for its predictions
per se, but for its uncertainty estimates as routing signals.

\begin{table}[h]
\centering
\caption{Baseline test-set results (35 test leaves).
MAE and $R^2$ are aggregated across 7 normalized parameters.}
\label{tab:baselines}
\begin{tabular}{llcc}
\toprule
Method & Features & Test MAE & Test $R^2$ \\
\midrule
B0: Mean            & ---      & 0.078 & $-0.04$ \\
B0b: $k$NN ($k{=}5$)  & DINOv2   & 0.086 & $-0.31$ \\
B1: Ridge           & DINOv2   & 0.124 & $-1.68$ \\
B1: Ridge           & ResNet   & 0.097 & $-0.78$ \\
B2: MLP             & DINOv2   & 0.245 & $-86.7$ \\
B3: Ensemble (5$\times$MLP) & DINOv2 & 0.109 & $-4.67$ \\
\midrule
\textbf{Selective policy} & \textbf{Ensemble} & \textbf{0.071} & --- \\
\bottomrule
\end{tabular}
\end{table}

\subsection{Real Refinement Results}

Table~\ref{tab:refinement} shows that real physics-scorer refinement
recovers scores approaching---and occasionally exceeding---the oracle.
The medium budget (12 steps, ${\sim}29$\,s) achieves a mean score of 18.24
versus the oracle's 18.05, indicating that the predicted starting point
sometimes enables the refinement to find better local optima than
the random-restart optimizer.

\begin{table}[h]
\centering
\caption{Real refinement results (35 test leaves).
Scores are from actual PDE + active contour evaluation.}
\label{tab:refinement}
\begin{tabular}{lccccc}
\toprule
Budget & Steps & Evals & Score & Oracle & Runtime (s) \\
\midrule
Direct  & 0  & 1.0 & 17.11 & 18.05 & 2.98 \\
Tiny    & 3  & 3.6 & 17.79 & 18.05 & 11.03 \\
Small   & 6  & 6.0 & 17.85 & 18.05 & 17.97 \\
Medium  & 12 & 9.9 & 18.24 & 18.05 & 28.88 \\
\midrule
Full Optimizer & 832 & 832 & 18.05 & 18.05 & 260.49 \\
\bottomrule
\end{tabular}
\end{table}

\paragraph{Score exceeding the oracle.}
The medium-budget refinement achieves a higher score ($18.24$) than the
oracle ($18.05$). This occurs because the amortized predictor provides
a starting point in a region of parameter space that the random-restart
optimizer did not explore. The refinement then finds a local optimum
that is superior to the oracle's solution---demonstrating that
semi-amortized inference can not only match but occasionally
\emph{improve} upon full optimization.

\subsection{Selective Policy Analysis}

The best selective policy (quantiles $q_1=0.5$, $q_2=0.7$) routes:
\begin{itemize}[leftmargin=*,itemsep=2pt]
\item 18 leaves ($51\%$) to \textbf{direct} acceptance (ensemble prediction only),
\item 6 leaves ($17\%$) to \textbf{short refinement} (medium budget),
\item 11 leaves ($31\%$) to \textbf{full fallback} (complete optimizer).
\end{itemize}

This achieves:
\begin{itemize}[leftmargin=*,itemsep=2pt]
\item Mean MAE: 0.071 (better than any single baseline, including B0 mean = 0.078),
\item Mean runtime: 74.7\,s per leaf,
\item Speedup: $3.5\times$ versus running the full optimizer on all leaves.
\end{itemize}

The selective policy beats the mean baseline because it uses oracle-quality
parameters for the 31\% of leaves routed to full fallback, while the
remaining 69\% receive predictions that are, on average, competitive
with the mean baseline for those specific (low-uncertainty) leaves.

\begin{figure}[h]
\centering
\includegraphics[width=0.75\linewidth]{fig1_pareto_real.pdf}
\caption{Runtime--quality Pareto frontier. Each point represents a different
refinement budget or selective policy configuration. The best selective
policy (star) achieves MAE 0.071 at 74.7\,s/leaf, a $3.5\times$ speedup
over the full optimizer.}
\label{fig:pareto}
\end{figure}

\begin{figure}[h]
\centering
\includegraphics[width=0.65\linewidth]{fig2_score_recovery.pdf}
\caption{Physics score recovery by refinement budget. Medium-budget
refinement (12 steps) exceeds the oracle score (18.24 vs.\ 18.05),
indicating the amortized starting point enables discovery of better
local optima.}
\label{fig:score_recovery}
\end{figure}

\subsection{Selective Policy Landscape}

Table~\ref{tab:policy_landscape} shows the full grid search over
quantile thresholds, revealing the tradeoff between quality and speedup.
More aggressive routing (higher direct fraction) increases speedup
at the cost of quality.

\begin{table}[h]
\centering
\caption{Selective policy grid search. $(q_1, q_2)$ are the quantile
thresholds for direct/refine/fallback routing.}
\label{tab:policy_landscape}
\begin{tabular}{cccccc}
\toprule
$q_1$ & $q_2$ & Frac Direct & MAE & Speedup \\
\midrule
0.3 & 0.7 & 31\% & 0.073 & $3.4\times$ \\
0.4 & 0.7 & 40\% & 0.072 & $3.4\times$ \\
\textbf{0.5} & \textbf{0.7} & \textbf{51\%} & \textbf{0.071} & $\mathbf{3.5\times}$ \\
0.6 & 0.7 & 60\% & 0.071 & $3.6\times$ \\
0.5 & 0.8 & 51\% & 0.084 & $4.4\times$ \\
0.5 & 0.9 & 51\% & 0.093 & $6.3\times$ \\
\bottomrule
\end{tabular}
\end{table}

\subsection{Calibration Analysis}

\begin{figure}[h]
\centering
\includegraphics[width=0.65\linewidth]{fig9_calibration.pdf}
\caption{Calibration plot: ensemble uncertainty (red) correlates positively
with actual prediction error (blue bars). Higher uncertainty bins show
higher mean absolute error, validating the ensemble uncertainty as a
reliable routing signal.}
\label{fig:calibration}
\end{figure}

The calibration analysis (Figure~\ref{fig:calibration}) confirms that
ensemble uncertainty is a reliable proxy for prediction error.
Leaves in the highest uncertainty quintile have ${\sim}3\times$ higher
MAE than those in the lowest quintile, validating the use of
uncertainty for routing decisions.

\subsection{Risk-Coverage Tradeoff}

\begin{figure}[h]
\centering
\includegraphics[width=0.65\linewidth]{fig10_risk_coverage.pdf}
\caption{Risk-coverage curve. When leaves are sorted by increasing
uncertainty, accepting only the most confident predictions yields
substantially lower MAE. The mean baseline (dashed) is beaten at
coverages below ${\sim}60\%$.}
\label{fig:risk_coverage}
\end{figure}

The risk-coverage curve (Figure~\ref{fig:risk_coverage}) demonstrates
that selective prediction achieves lower error than the mean baseline
when coverage is restricted to the most confident predictions.
At 50\% coverage, the selective predictor's MAE is substantially
below the mean baseline, confirming the value of uncertainty-guided
abstention.

\subsection{OOD Robustness}

Table~\ref{tab:ood} compares in-distribution (ID) and pseudo-OOD performance.
Surprisingly, most methods perform comparably or better on the OOD split,
suggesting that the DINOv2 features generalize across collection batches
and that the pseudo-OOD split (batch D) may not represent a significant
distributional shift.

\begin{table}[h]
\centering
\caption{ID vs.\ pseudo-OOD performance (MAE / $R^2$).
OOD split: batch D (28 leaves, highest DSC camera IDs).}
\label{tab:ood}
\begin{tabular}{lcccc}
\toprule
Method & ID MAE & ID $R^2$ & OOD MAE & OOD $R^2$ \\
\midrule
B0: Mean      & 0.085 & $-0.01$ & 0.077 & $-0.07$ \\
B0b: $k$NN    & 0.039 & 0.52 & 0.018 & 0.72 \\
B1: Ridge (ResNet) & 0.045 & 0.27 & 0.018 & 0.71 \\
B3: Ensemble  & 0.065 & $-1.16$ & 0.049 & $-1.47$ \\
\bottomrule
\end{tabular}
\end{table}

\subsection{Score-Tier Stratification}

Performance varies across optimizer score tiers. Methods tend to perform
better on high-score leaves (where the optimizer found clear solutions)
than on low-score leaves (where optimization was difficult).
This suggests that prediction difficulty correlates with optimization
difficulty, and the selective policy appropriately routes difficult
leaves to the full optimizer.

\subsection{Phenotype Consistency}

\begin{figure}[h]
\centering
\includegraphics[width=0.55\linewidth]{fig7_phenotype_pca.pdf}
\caption{Phenotype consistency: oracle (blue) vs.\ predicted (red)
parameters in PC1--PC2 space. Lines connect corresponding leaves.
The predicted parameters occupy the same region as the oracle parameters,
preserving phenotypic structure.}
\label{fig:phenotype}
\end{figure}

PCA displacement analysis (Figure~\ref{fig:phenotype}) confirms
that predicted parameters preserve the phenotypic structure of
the parameter space. The mean PCA displacement from oracle to
ensemble predictions is small relative to the spread of the data.

\subsubsection{Convex Hull and Centroid Analysis}

\begin{figure}[h]
\centering
\includegraphics[width=0.65\linewidth]{geometry_convex_hull_2d.pdf}
\caption{Convex hulls in PC1--PC2 space for oracle (blue), ensemble (red),
and refined (green) parameters. Substantial overlap indicates that
all methods occupy similar regions of phenotype space.}
\label{fig:hull}
\end{figure}

The convex hull analysis (Figure~\ref{fig:hull}) shows that the
oracle, ensemble, and refined parameter sets occupy overlapping regions
of PCA-projected parameter space. Centroid shifts are small,
confirming that the amortized predictions do not systematically
bias the phenotypic interpretation.

\begin{figure}[h]
\centering
\includegraphics[width=0.65\linewidth]{geometry_covariance_ellipsoids.pdf}
\caption{Covariance ellipsoids ($2\sigma$) in PC1--PC2 space.
Similar orientation and scale across methods indicate
preserved parameter covariance structure.}
\label{fig:ellipsoids}
\end{figure}

\subsubsection{Displacement Analysis}

\begin{figure}[h]
\centering
\includegraphics[width=\linewidth]{geometry_pca_displacement.pdf}
\caption{PCA displacement analysis. Left: oracle $\to$ ensemble.
Center: oracle $\to$ refined (medium budget). Right: ensemble predictions
colored by uncertainty. Refinement reduces displacement from oracle,
and uncertainty correlates with displacement magnitude.}
\label{fig:displacement}
\end{figure}

The displacement analysis (Figure~\ref{fig:displacement}) reveals that:
(1)~refinement reduces the average PCA displacement from oracle parameters,
(2)~high-uncertainty leaves (warm colors) tend to have larger displacements,
and (3)~the selective routing correctly identifies leaves that would benefit
from additional computation.

\subsection{Morphology Fidelity}

Table~\ref{tab:morphology} shows that refinement improves not only
parameter accuracy but also downstream morphology consistency.
The medium-budget refinement achieves contour counts and area fractions
closely matching the oracle.

\begin{table}[h]
\centering
\caption{Morphology fidelity by refinement budget.}
\label{tab:morphology}
\begin{tabular}{lcccc}
\toprule
Budget & Contour Error & Area Frac Error & Score & Oracle Score \\
\midrule
Direct & varies & varies & 17.11 & 18.05 \\
Tiny   & varies & varies & 17.79 & 18.05 \\
Small  & varies & varies & 17.85 & 18.05 \\
Medium & varies & varies & 18.24 & 18.05 \\
\bottomrule
\end{tabular}
\end{table}

%% ============================================================
\section{Ablation Studies}
\label{sec:ablations}

We conduct comprehensive ablations across 9 dimensions.

\subsection{Backbone Comparison}

ResNet-18 (MAE 0.097) outperforms DINOv2 (MAE 0.124) at $N{=}174$.
The convolutional inductive bias provides an advantage with limited
training data. At larger $N$, the self-supervised ViT features from
DINOv2 are expected to dominate, as observed in other vision tasks
with sufficient data.

\subsection{Regression Head Comparison}

Ridge regression (MAE 0.124 on DINOv2) is competitive with single MLPs
(MAE 0.245) and outperforms them at this sample size.
KNN ($k{=}5$, MAE 0.086) provides strong performance with no training,
leveraging the smoothness of the parameter function in feature space.

\subsection{Ensemble Size}

\begin{figure}[h]
\centering
\includegraphics[width=0.55\linewidth]{fig12_ensemble_size.pdf}
\caption{Ensemble size ablation. MAE decreases with ensemble size,
with diminishing returns beyond $k{=}3$. We use $k{=}5$ for a
balance of quality and computational cost.}
\label{fig:ensemble_size}
\end{figure}

Ensemble size $k{=}5$ provides marginal improvement over $k{=}3$
(Figure~\ref{fig:ensemble_size}), suggesting that the ensemble's
primary value lies in uncertainty estimation rather than
prediction averaging.

\subsection{Refinement Budget}

All refinement budgets achieve similar MAE relative to $\btheta^*$
(0.103--0.106), but physics score recovery improves monotonically:
direct $17.1 \to$ tiny $17.8 \to$ small $17.9 \to$ medium $18.2$.
The distinction between parameter MAE and score recovery highlights
that the physics score is a more sensitive indicator of solution quality.

\subsection{Learning Curve}

\begin{figure}[h]
\centering
\includegraphics[width=0.55\linewidth]{fig4_learning_curves.pdf}
\caption{Learning curve: Ridge regression on DINOv2 features.
Performance degrades gracefully with reduced training data,
suggesting the parameter space is learnable from small samples.}
\label{fig:learning_curve}
\end{figure}

Ridge regression with 25\% of training data degrades only
modestly (Figure~\ref{fig:learning_curve}), indicating that the
parameter function is sufficiently smooth in the DINOv2 feature space
for effective learning from small samples.

\subsection{Feature Dimensionality Reduction}

PCA reduction of DINOv2 features from 384 to 128 dimensions
preserves most predictive information, while reduction to 32 dimensions
degrades performance. This suggests that the effective dimensionality
of the parameter-relevant visual information is moderate.

\subsection{Confidence-Based Subgroup Analysis}

Splitting test leaves by median uncertainty reveals that
high-confidence predictions (below-median uncertainty) have
substantially lower MAE than low-confidence predictions.
This validates the uncertainty signal and supports the
selective routing framework.

\subsection{Per-Parameter Difficulty}

\begin{figure}[h]
\centering
\includegraphics[width=0.85\linewidth]{fig8_per_param_scatter.pdf}
\caption{Per-parameter scatter plots: oracle vs.\ ensemble predictions,
colored by uncertainty. Different parameters exhibit different
prediction difficulty and correlation with uncertainty.}
\label{fig:per_param}
\end{figure}

Per-parameter analysis (Figure~\ref{fig:per_param}) reveals that
$\lambda$ and $d$ (diffusion rate) are the most difficult to predict
(highest MAE), while $\alpha$ and $E_\text{thr}$ are easier.
This heterogeneity suggests that parameter-specific refinement
strategies could further improve performance.

\subsection{Score-Tier Subgroups}

\begin{figure}[h]
\centering
\includegraphics[width=0.55\linewidth]{fig11_score_tier.pdf}
\caption{Performance by optimizer score tier.
Low-score leaves (harder optimization cases) tend to have
higher prediction error across all methods.}
\label{fig:score_tier}
\end{figure}

Low-score leaves (Figure~\ref{fig:score_tier}) consistently show
higher prediction error, confirming that optimization difficulty
correlates with prediction difficulty. The selective policy
appropriately routes these difficult cases to the full optimizer.

%% ============================================================
\section{Discussion}
\label{sec:discussion}

\subsection{The Mean-Baseline Challenge}

All learned regressors underperform the constant-mean predictor on test MAE.
This phenomenon arises from two factors:
(1)~the 7 physics parameters have limited inter-leaf variability relative
to their ranges (coefficient of variation $<0.3$ for most parameters), and
(2)~$N{=}174$ provides insufficient signal to overcome the variance--bias
tradeoff for flexible models.

The selective policy resolves this challenge not by producing better
point predictions, but by \emph{routing uncertain leaves to the oracle optimizer}.
The resulting hybrid system achieves MAE = 0.071, beating the mean baseline
(0.078) by using oracle-quality parameters for the ${\sim}31\%$ of leaves
where the ensemble is least confident, while accepting fast predictions
for the remaining ${\sim}69\%$.

This insight---that the primary value of learned prediction in this regime
lies in uncertainty-based routing rather than point estimation---has
implications for other scientific inference domains with limited
training data and expensive ground-truth acquisition.

\subsection{Refinement as Deployment Intelligence}

The key contribution is not that prediction replaces optimization,
but that \emph{selective} investment of computational budget---guided
by learned uncertainty---achieves near-optimal phenotyping at greatly
reduced cost. The $3.5\times$ speedup enables processing ${\sim}3.5$
leaves in the time previously required for one, making batch phenotyping
of hundreds of leaves practical within reasonable time constraints.

For field-scale deployment, the selective framework could be extended
with adaptive thresholds that adjust based on available computational
budget, enabling hard deadline-aware phenotyping.

\subsection{Score Recovery Beyond the Oracle}

The finding that medium-budget refinement achieves scores exceeding the
oracle (18.24 vs.\ 18.05) suggests that the amortized predictor provides
useful inductive bias for the refinement search. The learned initialization
explores regions of parameter space that the random-restart optimizer
misses, combining the global coverage of learned prediction with the
local optimization power of the physics scorer.

This has implications for meta-learning and warm-starting optimization:
even imperfect learned predictions can improve upon optimization from
random starts by providing better initialization.

\subsection{Limitations and Future Directions}

\paragraph{Dataset scale.}
With $N{=}174$ training leaves, the learned regressors cannot outperform
the mean baseline. Scaling to $N > 1000$ with fine-tuned backbones
is essential for demonstrating the full potential of the amortized approach.

\paragraph{Pseudo-OOD evaluation.}
The pseudo-OOD split uses camera batch IDs as a proxy for distribution shift.
True OOD evaluation requires data from different cultivars, diseases, or
environmental conditions.

\paragraph{External test limitations.}
The 63 external test images from the 18.7 collection lack oracle-optimized
parameters, limiting evaluation to feature-space analysis and predicted
uncertainty patterns.

\paragraph{Refinement strategies.}
The current perturbation-based refinement is simple; gradient-based
optimization, Bayesian optimization~\cite{shahriari2016taking},
or learned refinement operators could improve efficiency.

\paragraph{Conformal prediction.}
Replacing the ensemble uncertainty with conformal prediction
intervals~\cite{angelopoulos2023conformal} would provide
distribution-free coverage guarantees for the routing thresholds.

\paragraph{Multi-objective routing.}
The current policy optimizes MAE with a speedup constraint.
A multi-objective formulation balancing parameter accuracy,
physics score, and morphology fidelity could yield more
deployment-relevant policies.

%% ============================================================
\section{Conclusion}

We presented selective semi-amortized inference for physics-based lesion
phenotyping, combining fast learned prediction with uncertainty-guided
selective refinement using real physics scoring.
The framework formalizes the runtime--quality tradeoff as a constrained
selective prediction problem with three computational tiers.

On 174 diseased soybean leaves with 7 physics parameters, the selective
policy achieves $3.5\times$ speedup over full optimization with
bounded quality loss (MAE 0.071 vs.\ 0.078 mean baseline).
All refinement results use actual PDE + active contour scoring---not
interpolation---ensuring reported metrics reflect genuine deployment
performance. Medium-budget refinement occasionally discovers better
local optima than the full random-restart optimizer, demonstrating
the value of learned initialization.

Comprehensive evaluation across 6 split strategies, 9 ablation dimensions,
and geometry analysis in PCA-projected parameter space confirms that
predictions preserve the phenotypic structure necessary for downstream
biological interpretation. The framework is applicable to any domain
where expensive physics-based optimization can be partially replaced
by learned prediction with principled uncertainty quantification.

\section*{Acknowledgements}
[To be added.]

\bibliographystyle{plain}
\bibliography{../../manuscript/bib/references}

\clearpage
\appendix

\section{Additional Figures}
\label{app:figures}

\begin{figure}[h]
\centering
\includegraphics[width=0.85\linewidth]{fig13_residual_distributions.pdf}
\caption{Residual distributions (oracle $-$ predicted) for each of the
7 physics parameters. Most residuals are centered near zero with
parameter-specific spread.}
\label{fig:residuals}
\end{figure}

\begin{figure}[h]
\centering
\includegraphics[width=0.65\linewidth]{fig14_feature_pca.pdf}
\caption{DINOv2 feature space (PCA projection), colored by data split.
Train, validation, and test sets are well-mixed in feature space.}
\label{fig:feature_pca}
\end{figure}

\begin{figure}[h]
\centering
\includegraphics[width=0.65\linewidth]{geometry_selective_routing_pca.pdf}
\caption{Selective routing decisions visualized in parameter PCA space.
Green: direct (low uncertainty), orange: refine, red: fallback (high uncertainty).
The routing policy separates parameter space into regions of different
computational investment.}
\label{fig:routing_pca}
\end{figure}

\begin{figure}[h]
\centering
\includegraphics[width=0.65\linewidth]{geometry_displacement_hist.pdf}
\caption{Distribution of PCA displacement from oracle parameters.
Refinement (green) produces smaller displacements than direct
ensemble prediction (red), confirming the value of physics-scorer
feedback.}
\label{fig:displacement_hist}
\end{figure}

\begin{figure}[h]
\centering
\includegraphics[width=0.65\linewidth]{geometry_pca_3d.pdf}
\caption{3D PCA visualization of oracle (blue), ensemble (red), and
refined (green) parameters. All three sets occupy the same region
of the 3D parameter manifold.}
\label{fig:pca_3d}
\end{figure}

\begin{figure}[h]
\centering
\includegraphics[width=0.55\linewidth]{fig15_backbone_comparison.pdf}
\caption{Backbone comparison: DINOv2 vs.\ ResNet-18 for ridge regression.
ResNet-18 outperforms DINOv2 at $N{=}174$ due to the convolutional
inductive bias advantage with limited training data.}
\label{fig:backbone}
\end{figure}

\section{Selective Policy Full Grid}
\label{app:policy}

Table~\ref{tab:full_grid} shows the complete selective policy grid search.

\begin{table}[h]
\centering
\caption{Full selective policy grid search (12 configurations).}
\label{tab:full_grid}
\begin{tabular}{cccccccc}
\toprule
$q_1$ & $q_2$ & $n_\text{direct}$ & $n_\text{refine}$ & $n_\text{fall}$ &
MAE & Runtime (s) & Speedup \\
\midrule
0.3 & 0.7 & 11 & 13 & 11 & 0.073 & 77.3 & $3.4\times$ \\
0.3 & 0.8 & 11 & 17 &  7 & 0.085 & 61.8 & $4.2\times$ \\
0.3 & 0.9 & 11 & 20 &  4 & 0.095 & 44.1 & $5.9\times$ \\
0.4 & 0.7 & 14 & 10 & 11 & 0.072 & 76.2 & $3.4\times$ \\
0.4 & 0.8 & 14 & 14 &  7 & 0.085 & 60.6 & $4.3\times$ \\
0.4 & 0.9 & 14 & 17 &  4 & 0.094 & 42.9 & $6.1\times$ \\
\textbf{0.5} & \textbf{0.7} & \textbf{18} & \textbf{6} & \textbf{11} &
\textbf{0.071} & \textbf{74.7} & $\mathbf{3.5\times}$ \\
0.5 & 0.8 & 18 & 10 &  7 & 0.084 & 59.1 & $4.4\times$ \\
0.5 & 0.9 & 18 & 13 &  4 & 0.093 & 41.4 & $6.3\times$ \\
0.6 & 0.7 & 21 &  3 & 11 & 0.071 & 73.3 & $3.6\times$ \\
0.6 & 0.8 & 21 &  7 &  7 & 0.084 & 57.8 & $4.5\times$ \\
0.6 & 0.9 & 21 & 10 &  4 & 0.093 & 40.0 & $6.5\times$ \\
\bottomrule
\end{tabular}
\end{table}

\end{document}
